\documentclass[10pt]{article}
\usepackage[a4paper,margin=18mm]{geometry}
\usepackage{amsmath,amssymb,booktabs,array,enumitem}
\usepackage[T1]{fontenc}\usepackage{lmodern}
\usepackage{tikz}\usetikzlibrary{shapes.gates.logic.US,positioning,arrows.meta}
\usepackage[hidelinks]{hyperref}
\setlength{\parindent}{0pt}\setlength{\parskip}{4pt}
\setlength{\emergencystretch}{2em}
\newcommand{\VDD}{V_{DD}}
\newcommand{\mos}[4]{\begin{scope}[shift={(#1,#2)}]
\draw (0,.24)--(0,-.24);\draw (0,.20)--(.25,.20)--(.25,.5);
\draw (0,-.20)--(.25,-.20)--(.25,-.5);\draw (-.16,.23)--(-.16,-.23);
\ifnum#4=1\draw (-.26,0) circle (.06);\draw (-.75,0)--(-.32,0);\else\draw (-.75,0)--(-.16,0);\fi
\node[anchor=east,font=\scriptsize] at (-.76,0) {$#3$};\end{scope}}
\newcommand{\ground}[2]{\draw (#1,#2)--(#1,{#2-.12});\draw ({#1-.17},{#2-.12})--({#1+.17},{#2-.12});\draw ({#1-.11},{#2-.20})--({#1+.11},{#2-.20});}
\newcommand{\rail}[3]{\draw (#1,#2)--(#1,{#2+.14});\node[above,font=\scriptsize] at (#1,{#2+.14}) {$#3$};}

\usetikzlibrary{calc,patterns}
\newcommand{\hm}[4]{\begin{scope}[shift={(#1,#2)}]
\draw(-.24,0)--(.24,0);\draw(-.2,0)--(-.2,-.25)--(-.5,-.25);
\draw(.2,0)--(.2,-.25)--(.5,-.25);\draw(-.23,.16)--(.23,.16);
\ifnum#4=1\draw(0,.26)circle(.06);\draw(0,.32)--(0,.75);\else\draw(0,.16)--(0,.75);\fi
\node[font=\scriptsize,above right]at(.30,.05){$#3$};\end{scope}}
\newcommand{\tg}[3]{\draw(#1-.24,#2-.13)rectangle(#1+.24,#2+.13);\node[above,font=\tiny,inner sep=2pt]at(#1,#2+.13){$e=#3$};}
\begin{document}
\begin{center}{\Large\bfseries 2021 Midterm -- Corrected Solutions}\\[3pt]EE3032 Introduction to VLSI, National Central University\end{center}
The complete original paper has six questions. Derivations below use the photographed problem statements first. The three-page corrected student PDF and the seven answer/lecture photographs are comparison evidence, not an official numerical model. The archive assigns the paper to 2021; the paper itself has no printed year.

\textbf{1. Short answers (35 points)}
\begin{enumerate}[label=(\alph*),leftmargin=7mm,itemsep=6pt]
\item Coat the wafer with photoresist; align the mask and expose; develop the resist; etch the exposed target film/silicon through the resist pattern; strip the resist. Exposure/development makes the temporary resist mask, and etching transfers it to the target surface.
\item Insert a grounded shield between the wires, and increase their spacing (or shorten the parallel overlap). Both reduce mutual capacitive coupling. A repeater can lower victim impedance, but it does not remove the coupling capacitance.
\item SOI places a thin active silicon layer above a buried insulator (usually buried oxide) and a handle wafer. This isolates source/drain regions from the bulk substrate, reduces junction area and associated leakage/capacitance, and removes a conducting bulk path. Fully depleted thin-body SOI also improves electrostatic control. Gate tunneling and all other leakage are not automatically eliminated.
\item The course correction lists thermal oxidation and chemical vapor deposition (CVD) of oxide. Thermal growth consumes silicon and can use dry oxygen or wet steam; CVD deposits an oxide film rather than oxidizing the underlying silicon. The distinction matters when describing the physical process.
\item Static CMOS is ratio-less for correct steady-state Boolean levels and gives full rail-to-rail restored outputs. In its ideal stable state one network is off, so there is no direct DC supply-to-ground path; real leakage remains.
\item The intended two components are leakage-current mechanisms: reverse-biased source/drain-to-body junction leakage and subthreshold conduction through an off channel. Thus $P_{static}=V_{DD}I_{leak}$. The output load capacitance causes switching energy, not static DC power; a voltage source by itself is not a leakage mechanism. Modern devices can also have gate leakage, but that does not replace the two course-listed mechanisms.
\item PMOS passes a strong high level to the pull-up rail, while NMOS passes a strong low level to ground. Complementary pull-up/pull-down networks therefore restore both rails without the threshold loss of an NMOS-only high or a PMOS-only low.
\item Via-first and via-last. Via-middle is another integration sequence, but only two are requested.
\item For a series NMOS stack with its body at ground, the device nearest the output can have the largest elevated source voltage and hence largest $V_{SB}$. Its threshold rises by $\gamma(\sqrt{2\phi_F+V_{SB}}-\sqrt{2\phi_F})$. The analogous PMOS stack has its body at $V_{DD}$. This is a stack/source-to-body effect: being attached to the output drain alone does not create body effect, and a grounded-source inverter NMOS has $V_{SB}=0$.
\item During plasma processing, a long temporarily floating conductor can collect charge and drive a damaging voltage across a connected thin gate oxide. Add an antenna protection diode to provide a discharge path; or split/re-route the long connection with a jumper to a later metal layer so the early fabrication-stage antenna area is smaller. These are fabrication antenna fixes, rather than RF antenna measures.
\item Wafer-to-wafer (W2W), die-to-wafer (D2W), and die-to-die (D2D), as expressly listed on the photographed lecture-board correction. Face-to-face/back-to-back describe orientation; TSV and CMP describe other operations, so neither list is the requested three assembly groupings.
\item $C_L=C_g+C_d+C_w$: fan-out input gate capacitance, diffusion/junction capacitance on the driving output, and interconnect/routing capacitance.
\item Shorted-gate (SG) FinFET, with the two side gates tied together, and independent-gate (IG) FinFET, with independently controlled side gates.
\end{enumerate}
\newpage
\textbf{2. Transmission-gate 8:1 MUX and NAND (15 points)}

\textbf{(a)} Use a balanced tree of seven 2:1 MUX cells. Each cell has two transmission gates: the top path is enabled by $\bar S$, the bottom path by $S$. In the figure, each TG box is the parallel PMOS/NMOS pair shown in the key below. Its label gives the NMOS enable $e$; the PMOS gate is the complement $\bar e$ shown in the key.
\begin{center}\begin{tikzpicture}[x=.86cm,y=.88cm,font=\scriptsize]
\tg{1}{7.55}{\overline{S_0}}\draw(1.24,7.55)--(1.8,7.55)--(1.8,7);
\tg{1}{6.45}{S_0}\draw(1.24,6.45)--(1.8,6.45)--(1.8,7);
\draw(1.8,7)--(2.15,7);\fill(1.8,7)circle(.025);
\tg{1}{5.55}{\overline{S_0}}\draw(1.24,5.55)--(1.8,5.55)--(1.8,5);
\tg{1}{4.45}{S_0}\draw(1.24,4.45)--(1.8,4.45)--(1.8,5);
\draw(1.8,5)--(2.15,5);\fill(1.8,5)circle(.025);
\tg{1}{3.55}{\overline{S_0}}\draw(1.24,3.55)--(1.8,3.55)--(1.8,3);
\tg{1}{2.45}{S_0}\draw(1.24,2.45)--(1.8,2.45)--(1.8,3);
\draw(1.8,3)--(2.15,3);\fill(1.8,3)circle(.025);
\tg{1}{1.55}{\overline{S_0}}\draw(1.24,1.55)--(1.8,1.55)--(1.8,1);
\tg{1}{0.44999999999999996}{S_0}\draw(1.24,0.44999999999999996)--(1.8,0.44999999999999996)--(1.8,1);
\draw(1.8,1)--(2.15,1);\fill(1.8,1)circle(.025);
\tg{5}{6.55}{\overline{S_1}}\draw(5.24,6.55)--(5.8,6.55)--(5.8,6);
\tg{5}{5.45}{S_1}\draw(5.24,5.45)--(5.8,5.45)--(5.8,6);
\draw(5.8,6)--(6.15,6);\fill(5.8,6)circle(.025);
\tg{5}{2.55}{\overline{S_1}}\draw(5.24,2.55)--(5.8,2.55)--(5.8,2);
\tg{5}{1.45}{S_1}\draw(5.24,1.45)--(5.8,1.45)--(5.8,2);
\draw(5.8,2)--(6.15,2);\fill(5.8,2)circle(.025);
\tg{9}{4.55}{\overline{S_2}}\draw(9.24,4.55)--(9.8,4.55)--(9.8,4);
\tg{9}{3.45}{S_2}\draw(9.24,3.45)--(9.8,3.45)--(9.8,4);
\draw(9.8,4)--(10.15,4);\fill(9.8,4)circle(.025);
\draw(-.2,7.55)node[left]{$D_0$}--(.76,7.55);
\draw(-.2,6.45)node[left]{$D_1$}--(.76,6.45);
\draw(-.2,5.55)node[left]{$D_2$}--(.76,5.55);
\draw(-.2,4.45)node[left]{$D_3$}--(.76,4.45);
\draw(-.2,3.55)node[left]{$D_4$}--(.76,3.55);
\draw(-.2,2.45)node[left]{$D_5$}--(.76,2.45);
\draw(-.2,1.55)node[left]{$D_6$}--(.76,1.55);
\draw(-.2,0.45)node[left]{$D_7$}--(.76,0.45);
\draw(2.15,7)--(3,7)--(3,6.55)--(4.76,6.55);
\draw(2.15,5)--(3,5)--(3,5.45)--(4.76,5.45);
\draw(2.15,3)--(3,3)--(3,2.55)--(4.76,2.55);
\draw(2.15,1)--(3,1)--(3,1.45)--(4.76,1.45);
\draw(6.15,6)--(7,6)--(7,4.55)--(8.76,4.55);
\draw(6.15,2)--(7,2)--(7,3.45)--(8.76,3.45);
\draw(10.15,4)--(11,4)node[right]{$Y$};
\node[align=center]at(1,-.65){First stage\\$S_0$};\node[align=center]at(5,-.65){Second stage\\$S_1$};\node[align=center]at(9,-.65){Third stage\\$S_2$};
\end{tikzpicture}\end{center}
\begin{center}\begin{tikzpicture}[x=.75cm,y=.75cm,font=\scriptsize]
\hm{1.5}{2}{P}{1}\hm{1.5}{.5}{N}{0}
\draw(.25,1.4)--(.25,1.75)--(1,1.75);\draw(.25,1.4)--(.25,.25)--(1,.25);\draw(2,1.75)--(2.8,1.75)--(2.8,1.4);\draw(2,.25)--(2.8,.25)--(2.8,1.4);\draw(-.4,1.4)--(.25,1.4);\draw(2.8,1.4)--(3.4,1.4);\node[left]at(-.4,1.4){in};\node[right]at(3.4,1.4){out};\node[above]at(1.5,2.75){$\bar e$};\node[above]at(1.5,1.25){$e$};
\end{tikzpicture}\end{center}
Generate $\bar S_0$, $\bar S_1$, and $\bar S_2$ with three ordinary static CMOS inverters. The core has 14 TGs = 28 MOS, plus 6 MOS for these select inverters. No output restoration buffer is required to obtain full rails with a TG; a buffer can be added for drive. Every selected data path crosses three TGs. With $i=4S_2+2S_1+S_0$, the output is $Y=D_i$.

\textbf{(b)} Connect $S_2=A$, $S_1=B$, $S_0=C$, and choose
\[D_0=D_1=\cdots=D_6=1,\qquad D_7=0.\]
Then $F=1$ except when $ABC=111$, so $F=\overline{ABC}$. Constants mean physical $V_{DD}$/ground connections. The eight-input three-TG branch design in the students' work is also logically valid; the tree uses fewer switches and shares internal stages. It is important to connect the three selects, not just draw the data constants.
\newpage
\textbf{3. Static CMOS from the truth table (15 points)}

The original on-set is $\{3,5,7,9,10,11,13,14,15\}$ in binary $ABCD$ order. Grouping the ones gives
\[\boxed{F=AC+AD+BD+CD=AC+D(A+B+C).}\]
For example, $ABCD=1010$ has $F=1$, so a solution containing only the three $D$ terms is incomplete. Another valid factorization is $F=(A+B+C)(A+D)(C+D)$.

Use a complementary complex gate to produce $G=\bar F$, followed by a static inverter to restore the positive $F$ polarity. The NMOS pull-down of $G$ is $(A\text{ series }C)$ in parallel with $D\text{ series }(A\parallel B\parallel C)$. Its PMOS dual is $(A\parallel C)$ in series with $D\parallel(A\text{ series }B\text{ series }C)$.
\begin{center}\begin{tikzpicture}[x=.8cm,y=.82cm,font=\scriptsize]
\mos{0}{0.0}{A}{0}\node[right]at(0.45,0.45){$\mathrm{N1}$};\node[above]at(0.25,0.5){$G$};\node[below]at(0.25,-0.5){$n1$};
\mos{5}{0.0}{C}{0}\node[right]at(5.45,0.45){$\mathrm{N2}$};\node[above]at(5.25,0.5){$n1$};\node[below]at(5.25,-0.5){$0$};
\mos{10}{0.0}{D}{0}\node[right]at(10.45,0.45){$\mathrm{N3}$};\node[above]at(10.25,0.5){$G$};\node[below]at(10.25,-0.5){$n2$};
\mos{0}{-2.25}{A}{0}\node[right]at(0.45,-1.8){$\mathrm{N4}$};\node[above]at(0.25,-1.75){$n2$};\node[below]at(0.25,-2.75){$0$};
\mos{5}{-2.25}{B}{0}\node[right]at(5.45,-1.8){$\mathrm{N5}$};\node[above]at(5.25,-1.75){$n2$};\node[below]at(5.25,-2.75){$0$};
\mos{10}{-2.25}{C}{0}\node[right]at(10.45,-1.8){$\mathrm{N6}$};\node[above]at(10.25,-1.75){$n2$};\node[below]at(10.25,-2.75){$0$};
\mos{0}{-4.5}{A}{1}\node[right]at(0.45,-4.05){$\mathrm{P1}$};\node[above]at(0.25,-4.0){$V_{DD}$};\node[below]at(0.25,-5.0){$p1$};
\mos{5}{-4.5}{C}{1}\node[right]at(5.45,-4.05){$\mathrm{P2}$};\node[above]at(5.25,-4.0){$V_{DD}$};\node[below]at(5.25,-5.0){$p1$};
\mos{10}{-4.5}{D}{1}\node[right]at(10.45,-4.05){$\mathrm{P3}$};\node[above]at(10.25,-4.0){$p1$};\node[below]at(10.25,-5.0){$G$};
\mos{0}{-6.75}{A}{1}\node[right]at(0.45,-6.3){$\mathrm{P4}$};\node[above]at(0.25,-6.25){$p1$};\node[below]at(0.25,-7.25){$p2$};
\mos{5}{-6.75}{B}{1}\node[right]at(5.45,-6.3){$\mathrm{P5}$};\node[above]at(5.25,-6.25){$p2$};\node[below]at(5.25,-7.25){$p3$};
\mos{10}{-6.75}{C}{1}\node[right]at(10.45,-6.3){$\mathrm{P6}$};\node[above]at(10.25,-6.25){$p3$};\node[below]at(10.25,-7.25){$G$};
\mos{0}{-9.0}{G}{1}\node[right]at(0.45,-8.55){$\mathrm{P7}$};\node[above]at(0.25,-8.5){$V_{DD}$};\node[below]at(0.25,-9.5){$F$};
\mos{5}{-9.0}{G}{0}\node[right]at(5.45,-8.55){$\mathrm{N7}$};\node[above]at(5.25,-8.5){$F$};\node[below]at(5.25,-9.5){$0$};
\end{tikzpicture}\end{center}
All source/drain names in this schematic are electrical net labels: repeated names are connected, and $0$ denotes ground. This is a complete 14-MOS realization. The six NMOS implement $F$ as a conduction condition; the six PMOS implement $\bar F$. The last two transistors invert $G$. PMOS bodies connect to $V_{DD}$ and NMOS bodies to ground. Exhaustive connectivity checking gives exactly one rail connection and the correct $F$ for all 16 rows.

Both student sets have valid four-term Boolean expressions (one includes a redundant product). The corrected circuit is factored and includes the required output polarity explicitly; it does not assume complemented inputs are freely available.
\newpage
\textbf{4. Feedback NMOS gate (15 points: 9+6)}
\begin{center}\begin{tikzpicture}[x=.85cm,y=.85cm,font=\scriptsize]
\draw[dashed](.15,4.5)rectangle(1.65,5.4);\node[align=center]at(.9,4.95){P\\network};\draw(.9,5.4)--(.9,5.8);\node[above]at(.9,5.8){1.2 V};
\draw(.9,4.5)--(.9,4.1)--(1.25,4.1)--(1.25,3.5);\mos{1}{3}{N3}{0}
\draw(1.25,2.5)--(1.25,2);\mos{1}{1.5}{N1}{0}\ground{1.25}{1}
\draw(-.75,2.25)--(.05,2.25)--(.05,3)--(.25,3);\draw(.05,2.25)--(.05,1.5)--(.25,1.5);\fill(-.75,2.25)circle(.045);\node[left]at(-.75,2.25){$V_{in}$};
\draw(1.25,4.1)--(4.6,4.1);\fill(4.6,4.1)circle(.045);\node[right]at(4.6,4.1){$V_{out}$};
\hm{3.25}{2.25}{N2}{0}\draw(1.25,2.25)--(2.75,2);\fill(1.25,2.25)circle(.035);\draw(3.75,2)--(4.1,2)--(4.1,2.6);\node[above]at(4.1,2.6){1.2 V};\draw(3.25,3)--(3.25,4.1);\fill(3.25,4.1)circle(.03);\node at(2.2,.25){(a)};
\begin{scope}[shift={(7.2,1.6)},x=.85cm,y=1.7cm]
\draw[->](0,0)--(4.25,0)node[right]{$t$};\draw[->](0,0)--(0,1.65)node[above]{Voltage};
\draw[thick](0,1.2)--(1.65,1.2)--(2.85,0)--(4,0);\node at(.75,1.36){$V_{out}$};
\draw[thick](1.1,0)--(2.55,1.2)--(4,1.2);\node at(3.7,1.4){$V_{in}$};
\draw[dashed](0,.4)--(1.58,.4)--(1.58,0);\draw[dashed](0,.8)--(2.07,.8)--(2.07,0);\node[left]at(0,.4){0.4 V};\node[left]at(0,1.2){1.2 V};\node[below]at(1.58,0){$t_1$};\node[below]at(2.07,0){$t_2$};\node at(2.07,.9){$x$};\end{scope}
\node at(8.75,.25){(b)};
\end{tikzpicture}\end{center}
Name the junction of N1, N3, and N2 $u$. N2's gate is $V_{out}$, not $V_{in}$; its right terminal is tied to 1.2 V.

\textbf{(a) Course fixed-resistance answer.} The photographed lecture-board correction expressly uses the divider with N1 and N2 conducting while $V_{out}$ is still high:
\[u=\frac{R_{N1}}{R_{N1}+R_{N2}}(1.2\text{ V})=\frac{R}{R+2R}(1.2\text{ V})=0.4\text{ V}.\]
N3 starts conducting when $V_{GS3}=V_{in}-u$ reaches $V_{tn}=0.4$ V. Thus the threshold used for the qualitative point $x$ is
\[\boxed{V_{in}(x)\simeq0.4+0.4=0.8\text{ V}.}\]
This is a lumped on-resistance approximation. The first student's shortcut $V_{out}-V_t$ happens to give the same number, but it omits the specified resistance ratio and is not the correct divider derivation. A time-constant/ramp argument is also unnecessary: no numerical input slew or capacitance is given.

\textbf{(b) Region conditions.} For an NMOS with the lower-potential terminal used as its source, cutoff means $V_{GS}<V_t$; saturation means $V_{GS}\ge V_t$ and $V_{GD}\le V_t$; triode means $V_{GS}\ge V_t$ and $V_{GD}>V_t$. Applying the actual gate connections gives:
\begin{center}\begin{tabular}{lll}\toprule Device & $V_{GS}$ & $V_{GD}$\\\midrule N1 & $V_{in}$ & $V_{in}-u$\\N2 & $V_{out}-u$ & $V_{out}-1.2$\\N3 & $V_{in}-u$ & $V_{in}-V_{out}$\\\bottomrule\end{tabular}\end{center}
\newpage
\textbf{4(b), continued: distinguish onset from output discharge.}

The interval ends at the drawn input/output crossing $x$, so before that crossing $V_{out}\ge V_{in}$. In the conducting output-discharge portion, N3 is on, which means $V_{in}-u>0.4$ V. This directly proves the expected conducting-state answer:
\begin{center}\begin{tabular}{lll}\toprule Device & Region during discharge before $x$ & Reason\\\midrule N1 & triode (linear) & $V_{GD1}=V_{in}-u>0.4$\\N2 & saturation & $V_{GS2}\ge V_{in}-u>0.4$, $V_{GD2}\le0$\\N3 & saturation & $V_{GS3}>0.4$, $V_{GD3}=V_{in}-V_{out}\le0$\\\bottomrule\end{tabular}\end{center}
\[\boxed{\text{Conducting discharge: N1 linear, N2 saturation, N3 saturation.}}\]
The first student's N1 linear and N2 saturation classifications can apply here, but N3 linear is wrong: before the crossing, its $V_{GD}$ is nonpositive, so an on N3 is saturated.

There is also an onset portion if $t_1$ is interpreted as the instant $V_{in}$ first reaches 0.4 V. Before N3 turns on, $V_{in}-u<0.4$ V. An on N1 is then saturated, N2 is saturated if on, and N3 is off. At $V_{in}-u=0.4$ V, N1 is at its saturation/triode boundary and N3 is at its turn-on boundary. These boundary states are not the fully conducting discharge state above. The original qualitative waveform does not specify $u(t)$ or locate N3 turn-on numerically relative to every part of the ramp; the region inequalities handle both phases without inventing a transient.

\textbf{Physical-model limit.} Constant $R$, $2R$, and $R$ do not constitute a self-consistent saturation-current model. The exact moving $u(t)$ and the exact crossing $V_{out}=V_{in}$ depend on the pull-up network, MOS current laws, body effect, capacitances, and slew, which the original does not supply. Hence 0.8 V is the documented course approximation, rather than an exact transistor-level transient prediction. The region inequalities remain the authoritative test if a more detailed waveform is provided. For instance, a square-law saturation model with $\beta_1/\beta_2=2$ would give a different initial N3 turn-on value, $0.4+0.8/(1+\sqrt2)\simeq0.7314$ V; that additional model is not specified by the exam and is not substituted for its requested resistor answer.

\textbf{5. Four-transistor circuit (10 points)}
\begin{center}\begin{tikzpicture}[x=.9cm,y=.9cm,font=\small]
\mos{3}{3}{\phantom{B}}{1}\mos{3}{1.5}{\phantom{B}}{0}
\draw(3.25,3.5)--(3.25,4)--(.15,4)node[left]{$A$};\draw(3.25,1)--(3.25,.5)--(.15,.5)node[left]{$\bar A$};
\draw(3.25,2.5)--(3.25,2)--(4,2)node[right]{$y$};\draw(3.25,2)--(3.25,2);\draw(3.25,2)--(3.25,2);\draw(3.25,2)--(3.25,2);\draw(3.25,2)--(3.25,2);
\draw(.15,2)--(1.28,2);\node[left]at(.15,2){$B$};\draw(1.28,1.76)rectangle(1.78,2.24);\draw(1.78,2)--(3.25,2);
\draw(1.53,2.24)--(1.53,2.65);\draw(1.53,2.76)circle(.055);\draw(1.53,2.82)--(1.53,4);\draw(1.53,1.76)--(1.53,.5);
\draw(.65,2)--(.65,1.25)--(2.1,1.25)--(2.1,3)--(2.25,3);\draw(2.1,1.5)--(2.25,1.5);\fill(.65,2)circle(.03);\fill(2.1,1.5)circle(.03);\fill(3.25,2)circle(.03);
\end{tikzpicture}\end{center}
When $A=0$, the transmission gate is on and passes $B$ to $y$. When $A=1$, the TG is off and the right-hand pair has conventional rails $A=1$, $\bar A=0$, giving $y=\bar B$. Consequently
\[\boxed{y=\bar A B+A\bar B=A\oplus B.}\]
The first student called it an inverter and received only partial credit; that conclusion ignores the $A=0$ pass mode. The second student's XOR conclusion is correct.

The pass mode has finite TG resistance and relies on the upstream $B$ driver; it provides no independent restoring gain or fixed drive strength. The switched supply inputs $A,\bar A$ must themselves be well-driven complementary signals. During imperfectly matched transitions, contention or a floating output can occur. Supply/body connections and fan-out therefore matter; this is not interchangeable with a conventional always-powered static XOR under arbitrary loading and slew. Complement generation for $\bar A$ also costs an inverter if it is not already present.
\newpage
\textbf{6. Positive function and transition-specific inverter (10 points)}

\textbf{(a)} Factor the specified positive function:
\[F=X(Y+Z)+XW=X(Y+Z+W).\]
Implement $G=\overline{X(Y+Z+W)}$ with an 8-MOS complementary gate, then use a 2-MOS inverter. This needs only the original inputs:
\begin{center}\begin{tikzpicture}[x=.8cm,y=.82cm,font=\scriptsize]
\mos{0}{0.0}{X}{0}\node[right]at(0.45,0.45){$\mathrm{N1}$};\node[above]at(0.25,0.5){$G$};\node[below]at(0.25,-0.5){$n1$};
\mos{5}{0.0}{Y}{0}\node[right]at(5.45,0.45){$\mathrm{N2}$};\node[above]at(5.25,0.5){$n1$};\node[below]at(5.25,-0.5){$0$};
\mos{10}{0.0}{Z}{0}\node[right]at(10.45,0.45){$\mathrm{N3}$};\node[above]at(10.25,0.5){$n1$};\node[below]at(10.25,-0.5){$0$};
\mos{0}{-2.25}{W}{0}\node[right]at(0.45,-1.8){$\mathrm{N4}$};\node[above]at(0.25,-1.75){$n1$};\node[below]at(0.25,-2.75){$0$};
\mos{5}{-2.25}{X}{1}\node[right]at(5.45,-1.8){$\mathrm{P1}$};\node[above]at(5.25,-1.75){$V_{DD}$};\node[below]at(5.25,-2.75){$G$};
\mos{10}{-2.25}{Y}{1}\node[right]at(10.45,-1.8){$\mathrm{P2}$};\node[above]at(10.25,-1.75){$V_{DD}$};\node[below]at(10.25,-2.75){$p1$};
\mos{0}{-4.5}{Z}{1}\node[right]at(0.45,-4.05){$\mathrm{P3}$};\node[above]at(0.25,-4.0){$p1$};\node[below]at(0.25,-5.0){$p2$};
\mos{5}{-4.5}{W}{1}\node[right]at(5.45,-4.05){$\mathrm{P4}$};\node[above]at(5.25,-4.0){$p2$};\node[below]at(5.25,-5.0){$G$};
\mos{10}{-4.5}{G}{1}\node[right]at(10.45,-4.05){$\mathrm{P5}$};\node[above]at(10.25,-4.0){$V_{DD}$};\node[below]at(10.25,-5.0){$F$};
\mos{0}{-6.75}{G}{0}\node[right]at(0.45,-6.3){$\mathrm{N5}$};\node[above]at(0.25,-6.25){$F$};\node[below]at(0.25,-7.25){$0$};
\end{tikzpicture}\end{center}
The NMOS network is $X$ in series with $Y\parallel Z\parallel W$. Its PMOS dual is $X$ in parallel with $Y$--$Z$--$W$ in series. The final inverter is essential to make $F$ positive. The students' complemented-input versions can also be correct if the complements are explicitly generated, but expanding the unfactored expression needlessly duplicates $X$.

\textbf{(b)} The given function here is different: $F=\overline{X(Y+Z)}$. Its six-transistor AOI21 network is
\begin{center}\begin{tikzpicture}[x=.8cm,y=.82cm,font=\scriptsize]
\mos{0}{0.0}{X}{0}\node[right]at(0.45,0.45){$\mathrm{N_X}$};\node[above]at(0.25,0.5){$F$};\node[below]at(0.25,-0.5){$n$};
\mos{5}{0.0}{Y}{0}\node[right]at(5.45,0.45){$\mathrm{N_Y}$};\node[above]at(5.25,0.5){$n$};\node[below]at(5.25,-0.5){$0$};
\mos{10}{0.0}{Z}{0}\node[right]at(10.45,0.45){$\mathrm{N_Z}$};\node[above]at(10.25,0.5){$n$};\node[below]at(10.25,-0.5){$0$};
\mos{0}{-2.25}{X}{1}\node[right]at(0.45,-1.8){$\mathrm{P_X}$};\node[above]at(0.25,-1.75){$V_{DD}$};\node[below]at(0.25,-2.75){$F$};
\mos{5}{-2.25}{Y}{1}\node[right]at(5.45,-1.8){$\mathrm{P_Y}$};\node[above]at(5.25,-1.75){$V_{DD}$};\node[below]at(5.25,-2.75){$p$};
\mos{10}{-2.25}{Z}{1}\node[right]at(10.45,-1.8){$\mathrm{P_Z}$};\node[above]at(10.25,-1.75){$p$};\node[below]at(10.25,-2.75){$F$};
\end{tikzpicture}\end{center}
Every device has $W/L$. For the specified sequence $101\to100\to101$, only $Z$ switches, while $X=1,Y=0$. The truth sequence is $F:0\to1\to0$, so the equivalent inverter input is $Z$.
\newpage
\textbf{6(b), continued.}

On the rising-output edge ($Z:1\to0$), $P_X$ is off and $P_Y,P_Z$ are in series. Thus $R_{p,eq}=2R_p$ and $(W/L)_{p,eq}=W/(2L)$.

On the falling-output edge ($Z:0\to1$), $N_Y$ is off and $N_X,N_Z$ are in series. Thus $R_{n,eq}=2R_n$ and $(W/L)_{n,eq}=W/(2L)$.
\[\boxed{(W/L)_{p,eq}=(W/L)_{n,eq}=\frac{W}{2L},\quad\text{input}=Z.}\]
\begin{center}\begin{tikzpicture}[x=.85cm,y=.8cm,font=\scriptsize]
\mos{1}{2.5}{\phantom{Z}}{1}\mos{1}{1}{\phantom{Z}}{0}
\draw(1.25,3)--(1.25,3.35);\node[above]at(1.25,3.35){$V_{DD}$};\ground{1.25}{.5}
\draw(1.25,2)--(1.25,1.5);\draw(1.25,1.75)--(3.75,1.75)node[right]{$F$};\fill(1.25,1.75)circle(.03);
\draw(-.7,1.75)node[left]{$Z$}--(.05,1.75);\draw(.05,1.75)--(.05,2.5)--(.25,2.5);\draw(.05,1.75)--(.05,1)--(.25,1);
\node[right]at(1.55,2.5){$W/(2L)$};\node[right]at(1.55,1){$W/(2L)$};
\draw(3.3,1.75)--(3.3,.8);\draw(3.05,.8)--(3.55,.8);\draw(3.05,.66)--(3.55,.66);\ground{3.3}{.66}\node[right]at(3.55,.73){$C_L$};
\end{tikzpicture}\end{center}
Both widths may be written $W/2$ with length $L$, or widths $W$ with length $2L$ within the simple resistance model. PMOS and NMOS do not thereby have equal physical resistance: their mobilities may differ. Retain the original output capacitance $C_L$ when using
\[t_r\simeq0.69(2R_p)C_L,\qquad t_f\simeq0.69(2R_n)C_L.\]
These expressions capture the effective output drive, not a full Elmore network including internal-node charge or nonlinear MOS transients. No numerical time is determined without a load and device model.

The three-page student PDF uses the all-input-on pull-down and all-input-off pull-up equivalents, $2W/(3L)$ and $3W/(2L)$. Those are general network states, not the specified $X=1,Y=0$ transition, and therefore answer the wrong question. The five-page answer/lecture set's page 4 correctly selects the two series devices on each edge.

\textbf{Verification and source limits}

\texttt{verify.py} independently checks all 16 Q3 truth rows against the printed table and transistor connectivity, all 16 Q6(a) input states against its transistor graph, all 8 Q6(b) states plus the three transition states, all four Q5 conduction states, and all eight NAND select rows. The 8:1 selection tree is exhaustively checked for every 8-bit data vector and select value (2048 cases). The Q4 divider and device inequalities are checked separately under their labeled approximation. These checks are ideal Boolean/connectivity and algebra checks, not HSPICE measurements.

Point labels are preserved: Q1--Q6 totals are $35+15+15+15+10+10=100$, and Q4(a)+(b) prints $9+6=15$.

\textbf{Source notes.} The supplied comparison evidence includes two corrected student sets and three lecture-board correction photographs. Lecture-board pages specifically support Q1(d),(f),(k), and the Q4 resistance derivation. Later typed study annotations on the student images are distinguished from grading and from printed exam text.
\end{document}
